\section{An axiom set that was not minimal}

The interpretation was first stated with four frame conditions. A check that
drops each condition and re-runs the proofs reported two of them inert: dropping
either lost no theorem --- where ``lost'' then meant ``no longer proved'', which is
the wrong criterion and is corrected below. Neither was a reprieve. Dropping \emph{both} loses three
theorems, because each was derivable from the other three, so the four were a
redundant presentation of the same constraint. The interpretation now rests on
three independent conditions --- every coordinate supports the certificate, every
edge runs from a coordinate to the certificate, and the certificate is not one
of its own coordinates --- and the fourth, that the certificate supports nothing,
is discharged as a theorem about them. Each of the three loses at least one
theorem when dropped, and which theorem each carries is recorded.

An axiom that no theorem needs is either decoration or a redundant presentation
of an axiom that is doing the work, and the two are indistinguishable from
inside. The check that told them apart is worth more here than the axiom it
removed.

The criterion was corrected in stages because each earlier version could turn a
fact about search into a claim about an axiom. Counting any theorem that stopped
being \emph{proved} credited an axiom when the solver returned \texttt{unknown}.
Requiring an actual countermodel removed that error, but an open model search
was unstable. Bounding discovery to at most four nodes did not remove the
instability: CI run 32927946106 reported the edge restriction only
intermittently, and run 32946736266 returned \texttt{unknown} for its known
one-node witness under full-suite load, cascading to four P6 test failures. Both
adverse runs remain failure provenance; neither showed the condition inert.

The current measurement no longer discovers a witness. It verifies one pinned,
fully specified exact-domain countermodel for each condition. Without coordinate
support, a two-node world has an unchanged non-coordinate certificate, one
changed coordinate, identity reachability and no edges; this refutes withdrawal
by any damage. Without the edge restriction, a one-node world has a
non-coordinate certificate with a self-edge; this refutes the derived claim that
the certificate supports nothing. Without certificate non-coordinateness, a
one-node world makes the changed certificate itself a coordinate with a
self-edge; because changed nodes are not reopened, this again refutes withdrawal.
All three tables use reflexive reachability and zero ranks. The solver checks the
remaining axioms, the exact domain, every table entry and the negated named
theorem. A certificate is credited only when every repeated check returns an
actual countermodel; \texttt{unknown} remains \texttt{CANNOT\_CHECK}, and an
invalid supplied model is a failed certificate rather than evidence of
inertness.

These certificates establish local logical necessity relative to this
formalization and theorem set only. They were written and checked in the same
lane. External or independent validation remains \texttt{CANNOT\_CHECK} under
P6-U-T4; the certificates provide neither empirical evidence nor external
scientific authority.
